% !TEX program = xelatex
% ============================================================================
%  POSN Computer Qualification — Equations Slide Deck
%  Topic: สมการและอสมการ (Equations & Inequalities)
% ============================================================================

\documentclass[11pt, aspectratio=169]{beamer}

% --- Load shared preamble ---
% ============================================================================
%  POSN Computer Qualification — Shared Beamer Preamble
%  Used by all topic slide decks. Compile with XeLaTeX.
% ============================================================================

% ---------- Thai language & font ----------
\usepackage{iftex}
\usepackage{fontspec}

% XeTeX-specific: enable Thai line breaking
\ifXeTeX
  \XeTeXlinebreaklocale{th}
\fi
% Noto Sans Thai — body text (clean modern sans-serif, handles all Thai)
\setmainfont{Noto Sans Thai}[
  UprightFont = Noto Sans Thai Light,
  BoldFont    = Noto Sans Thai Medium,
  Scale       = 0.9
]
\setsansfont{Noto Sans Thai}[
  UprightFont = Noto Sans Thai Light,
  BoldFont    = Noto Sans Thai Medium,
  Scale       = 0.9
]

% Noto Sans Thai — clean modern sans-serif for structural elements
\newfontfamily\headingfont{Noto Sans Thai}[
  UprightFont = Noto Sans Thai Medium,
  BoldFont    = Noto Sans Thai Bold,
  Scale       = 0.9
]

% --- Heading font for structural elements ---
\setbeamerfont{title}{family=\headingfont, series=\bfseries, size=\huge}
\setbeamerfont{subtitle}{family=\headingfont, series=\mdseries}
\setbeamerfont{author}{family=\headingfont, series=\mdseries}

% --- Custom Title Page ---
\defbeamertemplate{title page}{custom}[1][]{%
  \centering
  \vspace*{2cm}
  % Title — in white
  {\usebeamerfont{title}\color{white}\inserttitle\par}
  \vspace{0.6cm}
  % Subtitle — in white
  {\usebeamerfont{subtitle}\color{white}\insertsubtitle\par}
  \vspace{2.5cm}
  % Author — blue filled rounded box, white text
  \begin{center}
    \begin{tcolorbox}[arc=3mm, colback=boxBlue, colframe=boxBlue!80,
                      width=0.42\textwidth, halign=center,
                      left=1.5em, right=1.5em, top=0.8em, bottom=0.8em,
                      boxrule=1.5pt, coltext=white]
      \usebeamerfont{author}\insertauthor
    \end{tcolorbox}
  \end{center}
}
\setbeamertemplate{title page}[custom]
\setbeamerfont{frametitle}{family=\headingfont, series=\bfseries}
\setbeamerfont{section in toc}{family=\headingfont}
\setbeamerfont{page number in head/foot}{family=\headingfont}
\setbeamerfont{section title}{family=\headingfont, series=\bfseries}
\setbeamerfont{subsection title}{family=\headingfont}

% Monospace font (for code/verbatim)
\setmonofont{Latin Modern Mono}[Scale=0.9]

% ---------- Math ----------
\usepackage{amsmath, amssymb}

% ---------- Graphics ----------
\usepackage{graphicx}
\usepackage{tikz}

% ---------- String parsing (for section headers) ----------
\usepackage{xstring}

% ---------- Colored boxes ----------
\usepackage[skins, breakable, xparse]{tcolorbox}
% Note: we do NOT load the 'most' option — it predefines example/solution environments
% that would conflict with our custom boxes.

% --- Color definitions ---
\definecolor{boxBlue}{RGB}{41, 128, 185}
\definecolor{boxGreen}{RGB}{39, 174, 96}
\definecolor{boxOrange}{RGB}{230, 126, 34}
\definecolor{boxPurple}{RGB}{142, 68, 173}
\definecolor{boxBlueBg}{RGB}{235, 245, 255}
\definecolor{boxGreenBg}{RGB}{235, 255, 240}
\definecolor{boxOrangeBg}{RGB}{255, 245, 235}
\definecolor{boxPurpleBg}{RGB}{248, 240, 255}

% --- Explanation box (blue) — for definitions, concepts, notes ---
\newtcolorbox{explanation}[1][]{
  enhanced,
  breakable,
  arc         = 2mm,
  colback     = boxBlueBg,
  colframe    = boxBlue!50,
  title       = #1,
  fonttitle   = \bfseries\color{white},
  coltitle    = white,
  colbacktitle= boxBlue,
  attach boxed title to top left = {yshift=-2mm, xshift=2mm},
  boxed title style = {arc=2mm, size=small},
  before skip = 8pt,
  after skip  = 8pt
}

% --- Example box (green) — for worked examples ---
\newtcolorbox{exambox}[1][]{
  enhanced,
  breakable,
  arc         = 2mm,
  colback     = boxGreenBg,
  colframe    = boxGreen!50,
  title       = #1,
  fonttitle   = \bfseries\color{white},
  coltitle    = white,
  colbacktitle= boxGreen,
  attach boxed title to top left = {yshift=-2mm, xshift=2mm},
  boxed title style = {arc=2mm, size=small},
  before skip = 8pt,
  after skip  = 8pt
}

% --- Solution box (orange) — for solutions and answers ---
\newtcolorbox{solbox}[1][]{
  enhanced,
  breakable,
  arc         = 2mm,
  colback     = boxOrangeBg,
  colframe    = boxOrange!50,
  title       = #1,
  fonttitle   = \bfseries\color{white},
  coltitle    = white,
  colbacktitle= boxOrange,
  attach boxed title to top left = {yshift=-2mm, xshift=2mm},
  boxed title style = {arc=2mm, size=small},
  before skip = 8pt,
  after skip  = 8pt
}

% --- Intuition box (purple) — for plain-language explanations, analogies ---
\newtcolorbox{intuition}[1][]{
  enhanced,
  breakable,
  arc         = 2mm,
  colback     = boxPurpleBg,
  colframe    = boxPurple!50,
  title       = #1,
  fonttitle   = \bfseries\color{white},
  coltitle    = white,
  colbacktitle= boxPurple,
  attach boxed title to top left = {yshift=-2mm, xshift=2mm},
  boxed title style = {arc=2mm, size=small},
  before skip = 8pt,
  after skip  = 8pt
}

% ---------- Beamer theme ----------
\usetheme{Madrid}
\usecolortheme{default}

% --- Custom beamer colors (blue accent matching explanation boxes) ---
\setbeamercolor{structure}{fg=boxBlue}
\setbeamercolor{palette primary}{fg=white, bg=boxBlue}
\setbeamercolor{palette secondary}{fg=white, bg=boxBlue!80}
\setbeamercolor{palette tertiary}{fg=white, bg=boxBlue!60}
\setbeamercolor{block title}{fg=white, bg=boxBlue}
\setbeamercolor{block body}{fg=black, bg=boxBlueBg}

% Remove navigation symbols for clean look
\setbeamertemplate{navigation symbols}{}

% Note: Frame title prefix removed — \addtobeamertemplate conflicts with beamer internals

% --- Outline (Table of Contents) styling ---
\setbeamerfont{section in toc}{family=\headingfont, series=\mdseries}
\setbeamerfont{subsection in toc}{family=\headingfont, series=\mdseries}
\setbeamertemplate{section in toc}{%
  \leavevmode\leftskip=1.5em
  \textcolor{boxBlue}{\textbf{\large\inserttocsectionnumber.}}%
  \quad\inserttocsection\par
  \vspace{0.2em}
}
\setbeamertemplate{subsection in toc}{%
  \leavevmode\leftskip=4em
  \textcolor{boxBlue!60}{\small$\bullet$}%
  \quad\inserttocsubsection\par
  \vspace{0.1em}
}

% Section header slides — Thai in white, English in smaller light blue below
\AtBeginSection{
  \begin{frame}[plain]{}
    \vspace*{2cm}
    \begin{beamercolorbox}[sep=12pt, center, rounded=true, shadow=true]{palette primary}
      \IfSubStr{\secname}{(}{%
        \StrBefore{\secname}{(}[\sectionthai]%
        \StrBetween[1,1]{\secname}{(}{)}[\sectioneng]%
        \begin{tabular}{@{}c@{}}
          {\usebeamerfont{title}\sectionthai} \\[0.2em]
          {\usebeamerfont{subtitle}\color{boxBlue!60}\sectioneng}
        \end{tabular}%
      }{%
        {\usebeamerfont{title}\secname\par}%
      }
    \end{beamercolorbox}
    \vfill
  \end{frame}
}

% Subsection header slides — smaller, centered
\AtBeginSubsection{
  \begin{frame}[plain]{}
    \vfill
    \centering
    {\usebeamerfont{section title}\insertsubsectionhead\par}
    \vfill
  \end{frame}
}

% Minimal footer: page number only, right-aligned
\setbeamertemplate{footline}{
  \hfill\usebeamerfont{page number in head/foot}\insertframenumber{} / \inserttotalframenumber\hspace*{1em}\vspace*{0.5ex}
}

% ---------- Spacing & readability ----------
\linespread{1.2}

% ---------- Useful commands ----------

% Highlight important text in blue
\newcommand{\highlight}[1]{\textcolor{boxBlue}{\textbf{#1}}}

% Math set notation helpers
\newcommand{\R}{\mathbb{R}}
\newcommand{\N}{\mathbb{N}}
\newcommand{\Z}{\mathbb{Z}}
\newcommand{\Q}{\mathbb{Q}}

% ============================================================================
%  End of preamble
% ============================================================================


% --- Document info ---
\title[สมการและอสมการ]{สมการและอสมการ}
\subtitle{Equations \& Inequalities — สมการ, ระบบสมการ, อสมการ, ค่าสัมบูรณ์}
\author[None W. (Yuan)]{None W. (Yuan)}
\date{}

\begin{document}

% ===========================================================================
%  TITLE & OUTLINE
% ===========================================================================
\begin{frame}[plain]
  \titlepage
\end{frame}

\begin{frame}{สารบัญ (Outline)}
  \tableofcontents
\end{frame}

\begin{frame}{ก่อนเรียน (Prerequisites)}
  \begin{explanation}[สิ่งที่ควรรู้ก่อนเรียน]
    \begin{itemize}
      \item \textbf{จำนวนจริง} — การดำเนินการพื้นฐาน, สมบัติของจำนวนจริง, ค่าสัมบูรณ์
      \item \textbf{การแยกตัวประกอบ} — การแยกตัวประกอบพหุนาม (จากเรื่องพหุนาม)
      \item \textbf{เลขยกกำลัง} — $a^2$, $\sqrt{a}$ และสมบัติ
    \end{itemize}
  \end{explanation}

  \begin{intuition}[สมการคืออะไร?]
    \textbf{สมการ} คือประโยคที่บอกความเท่ากันของสองนิพจน์
    การแก้สมการคือการหาค่าตัวแปรที่ทำให้สมการเป็นจริง
  \end{intuition}
\end{frame}

% ===========================================================================
%  SECTION 1: สมการเชิงเส้นตัวแปรเดียว
% ===========================================================================
\section{สมการเชิงเส้นตัวแปรเดียว (Linear Equations)}

\begin{frame}{สมการเชิงเส้นตัวแปรเดียว}
  \begin{explanation}[รูปแบบทั่วไป]
    $ax + b = 0$ \quad โดยที่ $a \neq 0$

    \textbf{คำตอบ:} $x = -\dfrac{b}{a}$ \quad (มีคำตอบเดียวเสมอ)
  \end{explanation}

  \begin{exambox}
    แก้สมการ $3x - 7 = 2x + 5$
    \begin{align*}
      3x - 7 &= 2x + 5 \\
      3x - 2x &= 5 + 7 \\
      x &= 12
    \end{align*}
    \textbf{ตรวจสอบ:} $3(12) - 7 = 29$, \; $2(12) + 5 = 29$ \quad $\checkmark$
  \end{exambox}
\end{frame}

\begin{frame}{หลักการแก้สมการ}
  \begin{intuition}[หลักการสำคัญ]
    \begin{itemize}
      \item ย้ายข้างตัวแปรไปไว้ด้านหนึ่ง ตัวเลขไปไว้อีกด้านหนึ่ง
      \item \highlight{เครื่องหมายจะกลับเมื่อย้ายข้าง}
      \item ถ้ามีวงเล็บ ให้กระจายวงเล็บก่อน
      \item ถ้ามีเศษส่วน ให้คูณตลอดด้วย ค.ร.น. ของตัวส่วน
    \end{itemize}
  \end{intuition}
\end{frame}

\begin{frame}{ตัวอย่าง: สมการที่มีวงเล็บ}
  \begin{exambox}
    โจทย์: $2(x - 3) = 4x + 2$
    \begin{align*}
      2x - 6 &= 4x + 2 \\
      2x - 4x &= 2 + 6 \\
      -2x &= 8 \\
      x &= -4
    \end{align*}
    \textbf{ตรวจสอบ:} $2(-4-3) = -14$, \; $4(-4)+2 = -14$ \quad $\checkmark$
  \end{exambox}
\end{frame}

\begin{frame}{ตัวอย่าง: สมการที่มีเศษส่วน}
  \begin{exambox}
    โจทย์: $\dfrac{x}{2} + \dfrac{x}{3} = 5$ \newline
    คูณตลอดด้วย ค.ร.น. ของ 2 และ 3 = 6:
    \begin{align*}
      6\left(\frac{x}{2} + \frac{x}{3}\right) &= 6 \times 5 \\
      3x + 2x &= 30 \\
      5x &= 30 \quad \Rightarrow \quad x = 6
    \end{align*}
  \end{exambox}

  \begin{intuition}[เทคนิค]
    ถ้ามีเศษส่วน ให้คูณตลอดด้วย ค.ร.น. ของตัวส่วน เพื่อกำจัดเศษส่วนก่อน
  \end{intuition}
\end{frame}

% ===========================================================================
%  SECTION 2: สมการกำลังสอง
% ===========================================================================
\section{สมการกำลังสอง (Quadratic Equations)}

\begin{frame}{สมการกำลังสอง}
  \begin{explanation}[รูปแบบทั่วไป]
    $ax^2 + bx + c = 0$ \quad โดยที่ $a \neq 0$
    \textbf{วิธีแก้:}
    \begin{enumerate}
      \item \textbf{แยกตัวประกอบ} — ใช้เมื่อแยกได้ง่าย
      \item \textbf{ทำให้เป็นกำลังสองสมบูรณ์}
      \item \textbf{สูตรกำลังสอง} — ใช้ได้ทุกกรณี
    \end{enumerate}
  \end{explanation}

  \begin{explanation}[สูตรกำลังสอง (Quadratic Formula)]
    \[
      x = \frac{-b \pm \sqrt{b^2 - 4ac}}{2a}
    \]

    โดยที่ $b^2 - 4ac$ เรียกว่า \highlight{ดิสคริมิแนนต์ (Discriminant)}
  \end{explanation}
\end{frame}

\begin{frame}{ตัวอย่าง: การแยกตัวประกอบ}
  \begin{exambox}
    $x^2 + 5x + 6 = 0$
    หา 2 จำนวนที่คูณกันได้ 6 บวกกันได้ 5: \quad 2 และ 3
    \begin{align*}
      (x + 2)(x + 3) &= 0 \\
      x = -2 \quad &\text{หรือ} \quad x = -3
    \end{align*}
  \end{exambox}

  \begin{exambox}
    $x^2 - 9  0$
    ผลต่างกำลังสอง: \quad $(x - 3)(x + 3) = 0$

    $x = 3$ หรือ $x = -3$
  \end{exambox}

  \begin{intuition}[หลักการแยกตัวประกอบ]
    ถ้า $A \times B = 0$ แล้ว $A = 0$ \textbf{หรือ} $B = 0$
  \end{intuition}
\end{frame}

\begin{frame}{ตัวอย่าง: สูตรกำลังสอง}
  \begin{exambox}
    โจทย์: $2x^2 + 3x - 2 = 0$
    $a = 2$, $b = 3$, $c = -2$

    \textbf{ดิสคริมิแนนต์:} $b^2 - 4ac = 9 - 4(2)(-2) = 9 + 16 = 25$

    \textbf{สูตรกำลังสอง:}
    \[
      x = \frac{-3 \pm \sqrt{25}}{2 \times 2} = \frac{-3 \pm 5}{4}
    \]
    \[
      x = \frac{2}{4} = \frac{1}{2} \quad \text{หรือ} \quad x = \frac{-8}{4} = -2
    \]
  \end{exambox}
\end{frame}

\begin{frame}{ความหมายของดิสคริมิแนนต์}
  \begin{explanation}
    Discriminant: $\Delta = b^2 - 4ac$
    \begin{center}
      \begin{tabular}{c|c|c}
        $\Delta$ & จำนวนคำตอบ & ลักษณะ \\
        \hline
        $\Delta > 0$ & 2 คำตอบ (ต่างกัน) & กราฟตัดแกน $x$ 2 จุด \\
        $\Delta = 0$ & 1 คำตอบ (ซ้ำ) & กราฟสัมผัสแกน $x$ \\
        $\Delta < 0$ & ไม่มีคำตอบที่เป็นจำนวนจริง & กราฟไม่ตัดแกน $x$
      \end{tabular}
    \end{center}
  \end{explanation}

  \begin{exambox}[ตัวอย่าง]
    \begin{itemize}
      \item $x^2 - 4x + 4 = 0$: $\Delta = 0$ $\rightarrow$ $x = 2$ (คำตอบเดียว)
      \item $x^2 + x + 1 = 0$: $\Delta = -3 < 0$ $\rightarrow$ ไม่มีคำตอบใน $\R$
    \end{itemize}
  \end{exambox}
\end{frame}

\begin{frame}{ตัวอย่าง: ข้อสอบจริง}
  \begin{exambox}[ข้อสอบปี 65 ข้อ 23]
    คาใดคือผลบวกของค่า $a$ ทั้งหมดที่ทำให้สมการ $4x^2+(a+3)x+25=0$ มีคำตอบเดียว
  \end{exambox}

  \begin{solbox}
      มีคำตอบเดียวเมื่อ Discriminant: $\Delta = b^2 - 4ac=0$ 
      \begin{align*}
        (a+3)^2-4(4)(25)&=0 \\
        a^2+6x+9-400&=0 \\
        a^2+6x+391&=0 \\
        (a+23)(a-17)&=0
      \end{align*}
      $\therefore a=-23$ หรือ $a=17$ ผลรวมคือ $-6$ 
  \end{solbox}
\end{frame}


% ===========================================================================
%  SECTION 3: ระบบสมการ
% ===========================================================================
\section{ระบบสมการ (Systems of Equations)}

\begin{frame}{ระบบสมการเชิงเส้น 2 ตัวแปร}
  \begin{explanation}[รูปแบบ]
    \[
      \begin{cases}
        a_1x + b_1y = c_1 \\
        a_2x + b_2y = c_2
      \end{cases}
    \]

    \textbf{วิธีแก้:}
    \begin{enumerate}
      \item \textbf{กำจัดตัวแปร (Elimination)} — บวก/ลบสมการเพื่อกำจัดตัวแปรหนึ่ง
      \item \textbf{แทนค่า (Substitution)} — แก้ตัวแปรหนึ่งจากสมการหนึ่ง แล้วแทนในอีกสมการ
    \end{enumerate}
  \end{explanation}
\end{frame}

\begin{frame}{ตัวอย่าง: ระบบสมการ}
  \begin{exambox}[โจทย์]
    แก้ระบบสมการ:
    \[
      \begin{cases}
        2x + y = 7 \\
        x - y = 2
      \end{cases}
    \]
  \end{exambox}

  \begin{solbox}[วิธีที่ 1: กำจัดตัวแปร]
    บวกสมการทั้งสอง: $(2x + y) + (x - y) = 7 + 2$

    $3x = 9 \quad \Rightarrow \quad x = 3$

    แทน $x = 3$ ใน $x - y = 2$: \; $3 - y = 2 \quad \Rightarrow \quad y = 1$

    \textbf{คำตอบ:} $(x, y) = (3, 1)$
  \end{solbox}
\end{frame}

\begin{frame}{ตัวอย่าง: ระบบสมการ — แทนค่า}
  \begin{exambox}[โจทย์]
    แก้ระบบสมการ: \quad $y = 2x - 1$, \;\; $3x + 2y = 12$
  \end{exambox}

  \begin{solbox}[วิธีทำ]
    แทน $y = 2x - 1$ ลงใน $3x + 2y = 12$:
    \begin{align*}
      3x + 2(2x - 1) &= 12 \\
      3x + 4x - 2 &= 12 \\
      7x &= 14 \quad \Rightarrow \quad x = 2
    \end{align*}
    $y = 2(2) - 1 = 3$

    \textbf{คำตอบ:} $(x, y) = (2, 3)$
  \end{solbox}
\end{frame}

% ===========================================================================
%  SECTION 4: อสมการพื้นฐาน
% ===========================================================================
\section{อสมการ (Inequalities)}

\begin{frame}{อสมการคืออะไร?}
  \begin{explanation}[สมบัติของอสมการ]
    \begin{itemize}
      \item \textbf{บวก/ลบ ทั้งสองข้าง:} เครื่องหมายอสมการไม่เปลี่ยน
      \item \textbf{คูณ/หาร ด้วยจำนวนบวก:} เครื่องหมายอสมการไม่เปลี่ยน
      \item \highlight{คูณ/หาร ด้วยจำนวนลบ:} \textbf{เครื่องหมายอสมการกลับด้าน!}
    \end{itemize}
  \end{explanation}

  \begin{exambox}[ตัวอย่าง: การกลับเครื่องหมาย]
    \[
      -2x > 6 \quad \Rightarrow \quad x < -3
    \]

    (หารด้วย $-2$ ทั้งสองข้าง — เครื่องหมายกลับจาก $>$ เป็น $<$)
  \end{exambox}

  \begin{intuition}
    ถ้าคูณหรือหารด้วยจำนวนลบ \highlight{ต้องกลับเครื่องหมาย}
  \end{intuition}
\end{frame}

\begin{frame}{ตัวอย่าง: แก้อสมการเชิงเส้น}
  \begin{exambox}[โจทย์ 1: $3x - 5 < 2x + 4$]
    \begin{align*}
      3x - 2x &< 4 + 5 \\
      x &< 9
    \end{align*}
    \textbf{คำตอบ:} $x \in (-\infty, 9)$
  \end{exambox}
\end{frame}

\begin{frame}{ตัวอย่าง: อสมการที่ต้องกลับเครื่องหมาย}
  \begin{exambox}[โจทย์ 2: $5 - 2x \geq 3x + 10$]
    \begin{align*}
      -2x - 3x &\geq 10 - 5 \\
      -5x &\geq 5 \\
      x &\leq -1
    \end{align*}
    หารด้วย $-5$ — \highlight{กลับเครื่องหมายจาก $\geq$ เป็น $\leq$!}

    \textbf{คำตอบ:} $x \in (-\infty, -1]$
  \end{exambox}
\end{frame}

\begin{frame}{ตัวอย่าง: ข้อสอบจริง}
  \begin{exambox}[ข้อสอบปี 67 ข้อ 7]
    จากอสมการ $-43<-3x+17 \leq 5$ ถ้า $a$ และ $b$ คือ ค่ามากที่สุดและน้อยที่สุดที่เป็นจำนวนเต็มของ $x$ แล้วผลต่างของ $a$ และ $b$ มีค่าตรงกับข้อใด
  \end{exambox}

  \begin{solbox}
      $-43<-3x+17 \leq 5$\\
      $-60<-3x \leq -12$\\
      $20> x \geq 4$ หารด้วย $-3$ ต้องกลับเครื่องหมาย\\
      $\therefore x$ ที่เป็นจำนวนเต็มที่มากที่สุดคือ $4$ และน้อยที่สุดคือ $19$ ผลต่างคือ $15$ 
  \end{solbox}
\end{frame}

% ===========================================================================
%  SECTION 5: อสมการดีกรี > 1
% ===========================================================================
\section{อสมการดีกรีมากกว่า 1}

\begin{frame}{อสมการกำลังสอง: วิธีเส้นจำนวน (Sign Chart)}
  \begin{explanation}[ขั้นตอนการแก้อสมการกำลังสอง]
    \begin{enumerate}
      \item \textbf{จัดรูป}ให้ข้างหนึ่งเป็น 0
      \item \textbf{แยกตัวประกอบ}พหุนาม
      \item \textbf{หาจุดวิกฤต} (ค่า $x$ ที่ทำให้นิพจน์ = 0)
      \item \textbf{เขียนเส้นจำนวน} และทดสอบเครื่องหมายในแต่ละช่วง
      \item \textbf{เลือกช่วง}ที่ตรงตามเงื่อนไขอสมการ
    \end{enumerate}
  \end{explanation}

  \begin{intuition}[แนวคิด]
    พหุนามเปลี่ยนเครื่องหมายได้ที่จุดที่มันเท่ากับ 0 เท่านั้น
    ดังนั้นเราแบ่งเส้นจำนวนที่จุดวิกฤต แล้วทดสอบทีละช่วง
  \end{intuition}
\end{frame}

\begin{frame}{ตัวอย่าง: อสมการกำลังสอง $x^2 - x - 6 > 0$}
  \begin{exambox}[โจทย์]
    แก้อสมการ $x^2 - x - 6 > 0$
  \end{exambox}

  \begin{solbox}[วิธีทำ]
    \textbf{1. แยกตัวประกอบ:} $(x - 3)(x + 2) > 0$

    \textbf{2. จุดวิกฤต:} $x = -2$ และ $x = 3$ \quad (จุดที่นิพจน์ = 0)

    \textbf{3. ทดสอบเครื่องหมาย:}

    \small
    \begin{tabular}{c|ccc}
                     & $x<-2$ & $-2<x<3$ & $x>3$ \\
      \hline
      $(x-3)$        & $-$    & $-$      & $+$   \\
      $(x+2)$        & $-$    & $+$      & $+$   \\
      \hline
      ผลคูณ          & $+$    & $-$      & $+$
    \end{tabular}
    \normalsize

    \textbf{4. ต้องการ $> 0$:} เลือก $+$

    \textbf{คำตอบ:} $x \in (-\infty, -2) \cup (3, \infty)$
  \end{solbox}
\end{frame}

\begin{frame}{ตัวอย่าง: อสมการกำลังสอง $x^2 + 2x - 3 \leq 0$}
  \begin{exambox}[โจทย์]
    แก้อสมการ $x^2 + 2x - 3 \leq 0$
  \end{exambox}

  \begin{solbox}[วิธีทำ]
    \textbf{1. แยกตัวประกอบ:} $(x + 3)(x - 1) \leq 0$

    \textbf{2. จุดวิกฤต:} $x = -3$ และ $x = 1$

    \textbf{3. ทดสอบเครื่องหมาย:}

    \small
    \begin{tabular}{c|ccc}
                     & $x<-3$ & $-3<x<1$ & $x>1$ \\
      \hline
      $(x+3)$        & $-$    & $+$      & $+$   \\
      $(x-1)$        & $-$    & $-$      & $+$   \\
      \hline
      ผลคูณ          & $+$    & $-$      & $+$
    \end{tabular}
    \normalsize

    \textbf{4. ต้องการ $\leq 0$:} เลือก $-$ และ $= 0$

    \textbf{คำตอบ:} $x \in [-3, 1]$
  \end{solbox}
\end{frame}

\begin{frame}{ตัวอย่าง: อสมการดีกรี 3 — $x^3 - 4x \geq 0$}
  \begin{exambox}
    แก้อสมการ $x^3 - 4x \geq 0$
  \end{exambox}

  \begin{solbox}
    $x(x^2 - 4) \geq 0 \;\Rightarrow\; x(x-2)(x+2) \geq 0$

    \textbf{จุดวิกฤต:} $x = -2, 0, 2$

    \textbf{ทดสอบเครื่องหมาย (4 ช่วง):}

    \small
    \begin{tabular}{c|cccc}
                     & $x<-2$ & $-2<x<0$ & $0<x<2$ & $x>2$ \\
      \hline
      $(x+2)$        & $-$    & $+$      & $+$     & $+$   \\
      $x$            & $-$    & $-$      & $+$     & $+$   \\
      $(x-2)$        & $-$    & $-$      & $-$     & $+$   \\
      \hline
      ผลคูณ          & $-$    & $+$      & $-$     & $+$
    \end{tabular}
    \normalsize

    \textbf{ต้องการ $\geq 0$:} เลือก $+$ และจุดที่ $=0$

    \textbf{คำตอบ:} $x \in [-2, 0] \cup [2, \infty)$
  \end{solbox}
\end{frame}

\begin{frame}{ตัวอย่าง: ข้อสอบจริง}
  \begin{exambox}[ข้อสอบปี 65 ข้อ 43]
    จำนวนเต็มที่สอดคล้องกับ อสมการ $\dfrac{x(x-4)}{(x-1)(2x-1)} \leq 0$ มีทั้งหมดกี่จำนวน ($x$ เป็นจำนวนจริงใดๆ)
  \end{exambox}
  
\end{frame}

\begin{frame}{ตัวอย่าง: ข้อสอบจริง}
  \begin{solbox}[วิธีทำ]
      \textbf{จุดวิกฤต:} $x = 0, \dfrac{1}{2},1,4$
      \textbf{ทดสอบเครื่องหมาย (5 ช่วง):}

      \small
      \begin{tabular}{c|ccccc}
                      & $x<0$ & $0<x<\frac{1}{2}$ & $\frac{1}{2}<x<1$ & $1<x<4$ & $4<x$ \\
        \hline
        $x$          & $-$    & $+$      & $+$     & $+$     & $+$   \\
        $2x-1$       & $-$    & $-$      & $+$     & $+$     & $+$   \\
        $x-1$        & $-$    & $-$      & $-$     & $+$     & $+$   \\
        $x-4$        & $-$    & $-$      & $-$     & $-$     & $+$   \\
        \hline
        ผลคูณ          & $+$    & $-$      & $+$     & $-$     & $+$
      \end{tabular}
      \normalsize

      \textbf{ต้องการ $\leq 0$:} เลือก $-$
      $\therefore$ จำนวนเต็มที่สอดคล้อง ได้แก่ $0,2,3,4$ ทั้งหมด $4$ จำนวน 
      \\($\dfrac{1}{2},1$ ใช้ไม่ได้เนื่องจากจะทำให้ส่วนเป็น $0$) 
  \end{solbox}  
\end{frame}
% ===========================================================================
%  SECTION 6: ค่าสัมบูรณ์กับสมการ/อสมการ
% ===========================================================================
\section{ค่าสัมบูรณ์กับสมการและอสมการ}

\begin{frame}{สมการค่าสัมบูรณ์}
  \begin{explanation}[หลักการ]
    $|X| = a$ \quad ($a \geq 0$) \quad $\iff$ \quad $X = a$ หรือ $X = -a$

    \medskip
    ถ้า $a < 0$: \textbf{ไม่มีคำตอบ} (ค่าสัมบูรณ์ $\geq 0$ เสมอ)
  \end{explanation}

  \begin{exambox}
    ตัวอย่าง: $|2x - 5| = 7$
    \begin{align*}
      2x - 5 = 7 \quad &\Rightarrow \quad 2x = 12 \quad \Rightarrow \quad x = 6 \\
      2x - 5 = -7 \quad &\Rightarrow \quad 2x = -2 \quad \Rightarrow \quad x = -1
    \end{align*}
    \textbf{คำตอบ:} $x = 6$ หรือ $x = -1$
  \end{exambox}
\end{frame}

\begin{frame}{อสมการค่าสัมบูรณ์}
  \begin{explanation}[แบบที่ 1: $|X| < a$ ($a > 0$)]
    \[
      -a < X < a
    \]
    (ระยะห่างจาก 0 น้อยกว่า $a$ — อยู่ระหว่าง $-a$ กับ $a$)
  \end{explanation}

  \begin{explanation}[แบบที่ 2: $|X| > a$ ($a > 0$)]
    \[
      X < -a \quad \textbf{หรือ} \quad X > a
    \]
    (ระยะห่างจาก 0 มากกว่า $a$ — อยู่ข้างนอก)
  \end{explanation}
\end{frame}

\begin{frame}{ตัวอย่าง: อสมการค่าสัมบูรณ์}
  \begin{exambox}
    โจทย์ 1: $|x + 3| < 5$
    \begin{align*}
      -5 < x + 3 < 5 \\
      -8 < x < 2
    \end{align*}
    \textbf{คำตอบ:} $x \in (-8, 2)$
  \end{exambox}

  \begin{exambox}
    โจทย์ 2: $|2x - 1| \geq 3$
    \begin{align*}
      2x - 1 \leq -3 \quad &\Rightarrow \quad 2x \leq -2 \quad \Rightarrow \quad x \leq -1 \\
      2x - 1 \geq 3 \quad &\Rightarrow \quad 2x \geq 4 \quad \Rightarrow \quad x \geq 2
    \end{align*}
    \textbf{คำตอบ:} $x \in (-\infty, -1] \cup [2, \infty)$
  \end{exambox}
\end{frame}

\begin{frame}{ตัวอย่าง: ข้อสอบจริง}
  \begin{exambox}[ข้อสอบปี 67 ข้อ 16]
    ผลรวมของรากที่เป็นจำนวนเต็มของอสมการ $\lvert4x-5\rvert \leq 3x-2$ เท่ากับเท่าใด
  \end{exambox}

  \begin{solbox}
    \textbf{ราก} คือ คำตอบของสมการ
    \begin{align*}
      -3x+2 \leq 4x-5 \quad &\Rightarrow \quad 7 \leq 7x \quad \Rightarrow \quad 1 \leq x \\
      4x-5\leq 3x - 2 \quad &\Rightarrow \quad x \leq 3 \\
      1 \leq x \leq 3
    \end{align*}
    $\therefore$ รากที่เป็นจำนวนเต๋ช็ม คือ $1,2,3$ ซึ่งได้ผลรสมเป็น $6$
  \end{solbox}
\end{frame}
% ===========================================================================
%  SUMMARY
% ===========================================================================
\section{สรุป}

\begin{frame}{สรุปเนื้อหา}
  \begin{explanation}[สิ่งที่ได้เรียนรู้]
    \begin{enumerate}
      \item \textbf{สมการเชิงเส้น}: $ax + b = 0$ $\rightarrow$ $x = -b/a$
      \item \textbf{สมการกำลังสอง}: สูตร $x = \frac{-b \pm \sqrt{b^2-4ac}}{2a}$, ดิสคริมิแนนต์
      \item \textbf{ระบบสมการ}: กำจัดตัวแปร หรือแทนค่า
      \item \textbf{อสมการ}: คูณ/หารด้วยจำนวนลบ $\rightarrow$ \highlight{กลับเครื่องหมาย!}
      \item \textbf{อสมการดีกรี > 1}: \highlight{Sign Chart} — แยกตัวประกอบ, หาจุดวิกฤต, ทดสอบช่วง
      \item \textbf{ค่าสัมบูรณ์}: $|X| = a \iff X=\pm a$; \; $|X|<a \iff -a<X<a$
    \end{enumerate}
  \end{explanation}

  \begin{center}
    \large เอกสารนี้เป็นส่วนหนึ่งของเนื้อหาเตรียมสอบ \textbf{สอวน. คอมพิวเตอร์}
  \end{center}
\end{frame}

\end{document}
